\documentclass[lang=cn,newtx,10pt,sCheme=Chinese,citestyle=gb7714-2015, bibstyle=gb7714-2015]{../../Template/elegantbook}

% citestyle=gb7714-2015, bibstyle=gb7714-2015是为了将参考文献格式设置为国标GB7714-2015,不想使用这个文献格式标准的话可以删去

\title{概率论}
\subtitle{\,\,}

\author{实空}
\institute{无}
\date{\today}
\version{ElegantBook-4.5}
\bioinfo{全局显示/隐藏证明/解答/笔记}{\AllToggle}% 全局显示/隐藏证明/解答/笔记

\extrainfo{宠辱不惊,闲看庭前花开花落;
\\
去留无意,漫随天外云卷云舒.}


\setcounter{tocdepth}{3}

\logo{logo-blue.png}
\cover{cover.png}

% 修改标题页的橙色带
\definecolor{customcolor}{RGB}{32,178,170}
\colorlet{coverlinecolor}{customcolor}

%% 本文档额外使用的宏包和命令

\addbibresource[location=local]{reference.bib}% 引用参考文献

\usepackage{../../Styles/mystyle-elegantbook}

% 交叉引用设置（这里不要引用这个文件自身）
% 注意这里的引用地址不是.sty文件相对于其他被引用的.tex文件地址,而是需要引用的.tex文件相对于其他被引用的.tex文件地址!!!
\externaldocument[Abstract-Algebra-]{../Abstract-Algebra/main.aux}
\externaldocument[Analytic-Geometry-]{../Analytic-Geometry/main.aux}
\externaldocument[Basis-of-Algebra-]{../Basis-of-Algebra/main.aux}
\externaldocument[Basis-of-Analytics-]{../Basis-of-Analytics/main.aux}
\externaldocument[Complex-Analysis-]{../Complex-Analysis/main.aux}
\externaldocument[Differential-Geometry-]{../Differential-Geometry/main.aux}
\externaldocument[Functional-Analysis-]{../Functional-Analysis/main.aux}
\externaldocument[Numerical-Analysis-]{../Numerical-Analysis/main.aux}
% \externaldocument[Probability-Theory-]{../Probability-Theory/main.aux}
\externaldocument[Real-Analysis-]{../Real-Analysis/main.aux}
\externaldocument[Set-Theory-]{../Set-Theory/main.aux}
\externaldocument[Topology-]{../Topology/main.aux}

% 注意如果上述跨文档交叉引用代码报错，先检查是不是引用的.tex文件没有生成正确的.aux文件！！！

\begin{document}

\maketitle
\frontmatter

\tableofcontents

\mainmatter% 将行为改回预期版本，并重置页码

\chapter{大数定律和中心极限定理}

\section{大数定律}

\begin{definition}[(弱)大数定律]\label{definition:(弱)大数定律}
称随机变量序列$\left\{ X_n,n\geqslant 1 \right\} $服从\textbf{(弱)大数定律},若存在数列$\left\{a_n,n\geqslant  1\right\}$和$\left\{b_n,n\geqslant  1\right\}$,满足$b_n \to +\infty$,使得当$n\to +\infty$时,有
\begin{align*}
\frac{S_n-a_n}{b_n}\overset{P}{\rightarrow}0
\end{align*}
成立,其中$S_n=\sum\limits_{k=1}^n{X_k}$.
\end{definition}
\begin{note}
通常不作特别说明的时候,我们称随机变量序列$\{X_n,n\geqslant  1\}$服从\textbf{(弱)大数定律},都是指$a_n=\boldsymbol{E}S_n,b_n=n$时的规范形式.即下面的定义.
\end{note}

\begin{definition}[(弱)大数定律的规范形式]\label{definition:(弱)大数定律的规范形式}
称随机变量序列$\left\{ X_n,n\geqslant 1 \right\}$服从\textbf{(弱)大数定律},若当$n\to +\infty$时,有
\begin{align*}
\frac{S_n}{n}-\frac{\boldsymbol{E}S_n}{n}\overset{P}{\rightarrow}0
\end{align*}
成立,其中$S_n=\sum\limits_{k=1}^n{X_k}$.

即
\begin{align*}
\underset{n\rightarrow +\infty}{\lim}P\left( \left| \frac{1}{n}\sum_{k=1}^n{X_k}-\frac{1}{n}\boldsymbol{E}\sum_{k=1}^n{X_k} \right|<\varepsilon \right) =1.
\end{align*}
\end{definition}

\begin{theorem}[Bernoulli大数定律]\label{theorem:Bernoulli大数定律}
设在一次试验中事件$A$发生的概率为$p$,记$S_n$表示$n$次独立的这种试验$A$发生的次数,则当$n\to +\infty$时,
\begin{align*}
\frac{S_n}{n}\overset{P}{\rightarrow}p.
\end{align*}
\end{theorem}

\begin{theorem}[Markov大数定律]\label{theorem:Markov大数定律}
若随机变量序列$\left\{ X_n,n\geqslant 1 \right\}$满足$Markov$条件,即
\begin{align*}
\underset{n\rightarrow +\infty}{\lim}\frac{1}{n^2}\boldsymbol{D}\sum_{k=1}^n{X_k}=0.
\end{align*}
则$\left\{ X_n,n\geqslant 1 \right\}$服从(弱)大数定律.
\end{theorem}
\begin{proof}

对$\forall \varepsilon>0$,由Chebeshev不等式可知
\begin{align*}
1\geqslant P\left( \left| \frac{1}{n}\sum_{k=1}^n{X_k}-\frac{1}{n}\boldsymbol{E}\sum\limits_{k=1}^n{X_k} \right|<\varepsilon \right) \geqslant 1-\frac{\boldsymbol{D}\left( \frac{1}{n}\sum\limits_{k=1}^n{X_k} \right)}{\varepsilon ^2}=1-\frac{\boldsymbol{D}\sum\limits_{k=1}^n{X_k}}{n^2\varepsilon ^2}.
\end{align*}
令$n\to +\infty$,可得
\begin{align*}
\underset{n\rightarrow +\infty}{\lim}P\left( \left| \frac{1}{n}\sum_{k=1}^n{X_k}-\frac{1}{n}\boldsymbol{E}\sum_{k=1}^n{X_k} \right|<\varepsilon \right) =1.
\end{align*}

\end{proof}

\begin{theorem}[Chebeshev大数定律]\label{theorem:Chebeshev大数定律}
若随机变量序列$\left\{ X_n,n\geqslant 1 \right\}$两两不相关,且方差一致有界,
则$\left\{ X_n,n\geqslant 1 \right\}$服从(弱)大数定律.
\end{theorem}
\begin{proof}

因为随机变量序列$\left\{ X_n,n\geqslant 1 \right\}$方差一致有界,所以存在$C>0$,使得
\begin{align*}
\boldsymbol{D}X_k\leqslant C,k\in \mathbb{N}.
\end{align*}
又由于随机变量序列$\left\{ X_n,n\geqslant 1 \right\}$两两不相关,因此
\begin{align*}
\boldsymbol{D}\left( \frac{1}{n}\sum_{k=1}^n{X_k} \right) =\frac{1}{n^2}\boldsymbol{D}\sum_{k=1}^n{X_k}=\frac{1}{n^2}\sum_{k=1}^n{\boldsymbol{D}X_k}.
\end{align*}
根据$Chebeshev$不等式可得
\begin{align*}
1\geqslant P\left( \left| \frac{1}{n}\sum_{k=1}^n{X_k}-\frac{1}{n}\boldsymbol{E}\sum_{k=1}^n{X_k} \right|<\varepsilon \right) \geqslant 1-\frac{\boldsymbol{D}\left( \frac{1}{n}\sum\limits_{k=1}^n{X_k} \right)}{\varepsilon ^2}\geqslant 1-\frac{C}{n}.
\end{align*}
令$n\to +\infty$,可得
\begin{align*}
\underset{n\rightarrow +\infty}{\lim}P\left( \left| \frac{1}{n}\sum_{k=1}^n{X_k}-\frac{1}{n}\boldsymbol{E}\sum_{k=1}^n{X_k} \right|<\varepsilon \right) =1.
\end{align*}

\end{proof}

\begin{theorem}[Khinchin大数定律]\label{theorem:Khinchin大数定律}
若随机变量序列$\left\{ X_n,n\geqslant 1 \right\}$独立同分布,且数学期望存在,
则$\left\{ X_n,n\geqslant 1 \right\}$服从(弱)大数定律.
\end{theorem}
\begin{proof}



\end{proof}

\begin{theorem}[依概率收敛的性质]\label{theorem:依概率收敛的性质}
设\(\{ X_n, n \geqslant 1\}\)是概率空间\((\varOmega, \mathcal{F}, P)\)中的随机变量序列.

若\(X_n \overset{P}{\rightarrow} c\),且\(c \neq 0\),\(X_n \neq 0\),则\(\frac{1}{X_n} \overset{P}{\rightarrow} \frac{1}{c}\).
\end{theorem}
\begin{proof}

对\(\forall \varepsilon > 0\),由全概率公式可知
\begin{equation}\label{equation:1.1}
\begin{aligned}
&P\left(\left|\frac{1}{X_n} - \frac{1}{c}\right| \geqslant \varepsilon\right) = P\left(\left|\frac{X_n - c}{cX_n}\right| \geqslant \varepsilon\right) = P\left(\left|\frac{X_n - c}{cX_n - c^2 + c^2}\right| \geqslant \varepsilon\right) = P\left(\left|\frac{X_n - c}{c(X_n - c) + c^2}\right| \geqslant \varepsilon\right)
\\
&= P\left(\left|\frac{X_n - c}{c(X_n - c) + c^2}\right| \geqslant \varepsilon,\left|X_n - c\right| < \varepsilon\right) + P\left(\left|\frac{X_n - c}{c(X_n - c) + c^2}\right| \geqslant \varepsilon,\left|X_n - c\right| \geqslant \varepsilon\right)
\\
&\leqslant P\left(\left|\frac{X_n - c}{c(X_n - c) + c^2}\right| \geqslant \varepsilon,\left|X_n - c\right| < \varepsilon\right) + P\left(\left|X_n - c\right| \geqslant \varepsilon\right).
\end{aligned}
\end{equation}
当\(\left|\frac{X_n - c}{c(X_n - c) + c^2}\right| \geqslant \varepsilon\)且\(\left|X_n - c\right| < \varepsilon\)时,由\(\left|X_n - c\right| < \varepsilon\)可知\(-\varepsilon < X_n - c < \varepsilon\).
从而当\(c > 0\)时,\(c(X_n - c) > -c\varepsilon\);当\(c < 0\)时,\(c(X_n - c) > c\varepsilon\).因此对\(\forall c \neq 0\),都有\(c(X_n - c) > -|c|\varepsilon\).
于是此时\(\left|\frac{X_n - c}{c(X_n - c) + c^2}\right| < \left|\frac{X_n - c}{c^2 - |c|\varepsilon}\right|\).故
\begin{align*}
\left\{\left|\frac{X_n - c}{c(X_n - c) + c^2}\right| \geqslant \varepsilon,\left|X_n - c\right| < \varepsilon\right\} \subseteq \left\{\left|\frac{X_n - c}{c^2 - |c|\varepsilon}\right| \geqslant \varepsilon,\left|X_n - c\right| < \varepsilon\right\}.
\end{align*}
进而结合\eqref{equation:1.1}式,再根据概率的单调性可得
\begin{align*}
&P\left(\left|\frac{1}{X_n} - \frac{1}{c}\right| \geqslant \varepsilon\right) \leqslant P\left(\left|\frac{X_n - c}{c(X_n - c) + c^2}\right| \geqslant \varepsilon,\left|X_n - c\right| < \varepsilon\right) + P\left(\left|X_n - c\right| \geqslant \varepsilon\right)
\\
&\leqslant P\left(\left|\frac{X_n - c}{c^2 - |c|\varepsilon}\right| \geqslant \varepsilon,\left|X_n - c\right| < \varepsilon\right) + P\left(\left|X_n - c\right| \geqslant \varepsilon\right) \leqslant P\left(\left|\frac{X_n - c}{c^2 - |c|\varepsilon}\right| \geqslant \varepsilon\right) + P\left(\left|X_n - c\right| \geqslant \varepsilon\right)
\\
&= P\left(\left|X_n - c\right| \geqslant \varepsilon(c^2 - |c|\varepsilon)\right) + P\left(\left|X_n - c\right| \geqslant \varepsilon\right).
\end{align*}
令\(n \to +\infty\),则由\(\frac{1}{X_n} \overset{P}{\rightarrow} \frac{1}{c}\)可知
\begin{align*}
0 \leqslant P\left(\left|\frac{1}{X_n} - \frac{1}{c}\right| \geqslant \varepsilon\right) \leqslant P\left(\left|X_n - c\right| \geqslant \varepsilon(c^2 - |c|\varepsilon)\right) + P\left(\left|X_n - c\right| \geqslant \varepsilon\right) \to 0,n \to +\infty .
\end{align*}
故\(\lim_{n \to +\infty}P\left(\left|\frac{1}{X_n} - \frac{1}{c}\right| \geqslant \varepsilon\right) = 0\),即\(\frac{1}{X_n} \overset{P}{\rightarrow} \frac{1}{c}\).

\end{proof}

\begin{theorem}[依概率收敛与依分布收敛的关系]\label{theorem:依概率收敛与依分布收敛的关系}
设\(\{ X_n, n \geqslant 1\}\)和$X$分别是概率空间\((\varOmega, \mathcal{F}, P)\)中的随机变量序列和随机变量.

(1)$X_n\overset{P}{\rightarrow}X\text{蕴含}X_n\overset{d}{\rightarrow}X.$

\end{theorem}
\begin{proof}

(1)
记\(X_n\)(\(n\geqslant  1\))与\(X\)的分布函数分别为\(F_n(x)\)(\(n\geqslant  1\))与\(F\).
设\(X_n\overset{P}{\rightarrow}X\),往证\(F_n(x) \to F(x)\),\(n \to +\infty\),对所有\(F\)中连续点都成立.
为此只需证明对\(\forall x\in \mathbb{R}\),有
\begin{align*}
F(x - 0) \leqslant \varliminf_{n \to +\infty}F_n(x) \leqslant \varlimsup_{n \to +\infty}F_n(x) \leqslant F(x + 0).
\end{align*}
任取\(y < x\),由全概率公式可知
\begin{align*}
&F(y) = P(X\leqslant y) = P(X\leqslant y, X_n\leqslant x) + P(X\leqslant y, X_n > x)
\\
&\leqslant P(X_n\leqslant x) + P(|X_n - X|\geqslant x - y)
\\
&= F_n(x) + P(|X_n - X|\geqslant x - y)
\end{align*}
令\(n \to +\infty\),由\(X_n\overset{P}{\rightarrow}X\)可得\(F(y) \leqslant \varliminf_{n \to +\infty}F_n(x)\),对\(\forall y < x\)都成立.由于分布函数左右极限都存在,因此再令\(y \to x^-\),得到\(F(x - 0) \leqslant \varliminf_{n \to +\infty}F_n(x)\).

同理,任取\(y > x\),由全概率公式可知
\begin{align*}
&F_n(x) = P(X_n\leqslant x) = P(X_n\leqslant x, X\leqslant y) + P(X_n\leqslant x, X > y)
\\
&\leqslant P(X\leqslant y) + P(|X_n - X|\geqslant y - x)
\\
&= F(y) + P(|X_n - X|\geqslant y - x)
\end{align*}
令\(n \to +\infty\),由\(X_n\overset{P}{\rightarrow}X\)可得\(\varlimsup_{n \to +\infty}F_n(x) \leqslant F(y)\),对\(\forall y < x\)都成立.由于分布函数左右极限都存在,因此再令\(y \to x^+\),得到\(\varlimsup_{n \to +\infty}F_n(x) \leqslant F(x + 0)\).

又显然有\(\varliminf_{n \to +\infty}F_n(x) \leqslant \varlimsup_{n \to +\infty}F_n(x)\).
综上所述,我们有
\begin{align*}
F(x - 0) \leqslant \varliminf_{n \to +\infty}F_n(x) \leqslant \varlimsup_{n \to +\infty}F_n(x) \leqslant F(x + 0),\forall x\in \mathbb{R} .
\end{align*}
故结论得证.

\end{proof}

\begin{theorem}[特征函数性质]\label{theorem:特征函数性质}
特征函数在$(-\infty,+\infty)$上一致连续.
\end{theorem}
\begin{proof}

设随机变量\(X\)的概率密度函数为\(p(x)\),则\(\int_{-\infty}^{+\infty}p(x)\mathrm{d}x = 1\).其特征函数为\(f(t) = \int_{-\infty}^{+\infty}e^{itx}p(x)\mathrm{d}x\).

由反常积分收敛的柯西收敛准则,可知\(\forall \varepsilon > 0\),存在\(A > 0\),使得\(\int_{|x| > A}p(x)\mathrm{d}x < \frac{\varepsilon}{4}\).

于是对\(\forall t \in (-\infty, +\infty)\),取\(h = \frac{\varepsilon}{4A}\),我们有
\begin{equation}\label{equation:1.1123}
\begin{aligned}
|f(t + h) - f(t)| &= \left|\int_{-\infty}^{+\infty}(e^{i(t + h)x} - e^{itx})p(x)\mathrm{d}x\right|
= \left|\int_{-\infty}^{+\infty}e^{itx}(e^{ihx} - 1)p(x)\mathrm{d}x\right|\\
&\leqslant \int_{-\infty}^{+\infty}|e^{itx}||(e^{ihx} - 1)p(x)|\mathrm{d}x
= \int_{-\infty}^{+\infty}|(e^{ihx} - 1)|p(x)\mathrm{d}x
\end{aligned}
\end{equation}
又由欧拉公式,可得
\begin{equation}\label{equation:1.23434}
\begin{aligned}
|(e^{ihx} - 1)| &= |\cos(hx) - 1 + i\sin(hx)|
= \sqrt{(\cos(hx) - 1)^2 + \sin^2(hx)}\\
&= \sqrt{2 - 2\cos(hx)}
= \sqrt{4\sin^2\frac{hx}{2}} \leqslant 2\left|\sin\frac{hx}{2}\right|
\end{aligned}
\end{equation}
从而由\eqref{equation:1.1123}\eqref{equation:1.23434}式,可得
\begin{align*}
|f(t + h) - f(t)| &\leqslant \int_{-\infty}^{+\infty}|(e^{ihx} - 1)|p(x)\mathrm{d}x
\leqslant 2\int_{-\infty}^{+\infty}\left|\sin\frac{hx}{2}\right|p(x)\mathrm{d}x\\
&= 2\int_{|x| > A}\left|\sin\frac{hx}{2}\right|p(x)\mathrm{d}x + 2\int_{-A}^{A}\left|\sin\frac{hx}{2}\right|p(x)\mathrm{d}x\\
&\leqslant 2\int_{|x| > A}p(x)\mathrm{d}x + 2\int_{-A}^{A}|hx|p(x)\mathrm{d}x
\leqslant \frac{\varepsilon}{2} + 2hA\int_{-A}^{A}p(x)\mathrm{d}x\\
&\leqslant \frac{\varepsilon}{2} + 2hA\int_{-\infty}^{+\infty}p(x)\mathrm{d}x
= \frac{\varepsilon}{2} + 2hA < \varepsilon
\end{align*}
故特征函数\(f(t)\)在\((-\infty, +\infty)\)上一致连续. 

\end{proof}





\begin{lemma}[二次型期望公式]\label{lemma:二次型期望公式}
设 $\varepsilon$ 为 $n$ 维随机向量，$E(\varepsilon)=\mu$，$\operatorname{Cov}(\varepsilon)=\Sigma$，$A$ 为 $n\times n$ 实矩阵，则
\[
E(\varepsilon^T A \varepsilon) = \operatorname{tr}(A\Sigma) + \mu' A \mu.
\]
特别地，若 $\mu=0$，$\Sigma=\sigma^2 I_n$，则 $E(\varepsilon^TA\varepsilon)=\sigma^2\operatorname{tr}(A)$。
\end{lemma}
\begin{proof}

设$\varepsilon =\left( \varepsilon _1,\cdots ,\varepsilon _n \right) ^T,\mu =\left( \mu _1,\cdots ,\mu _n \right) ^T,\Sigma =\left( \Sigma _{ij} \right) _{n\times n},A=\left( a_{ij} \right) _{n\times n}$,则
\[
\varepsilon^T A \varepsilon = \sum_{i=1}^n\sum_{j=1}^n a_{ij}\varepsilon_i\varepsilon_j.
\]
从而
\begin{align}\label{eq::HJIOJW392234--32}
E(\varepsilon^T A\varepsilon )=\sum_{i,j=1}^n{a_{ij}E(\varepsilon _i\varepsilon _j)}.
\end{align}
由协方差定义知
\begin{align*}
\mathrm{Cov(}\varepsilon _i,\varepsilon _j)&=E\left[ \left( \varepsilon _i-E\left( \varepsilon _i \right) \right) \left( \varepsilon _j-E\left( \varepsilon _j \right) \right) \right] =E\left[ \left( \varepsilon _i-\mu _i \right) \left( \varepsilon _j-\mu _j \right) \right] 
\\
&=E\left( \varepsilon _i\varepsilon _j \right) -\mu _jE\left( \varepsilon _i \right) -\mu _iE\left( \varepsilon _j \right) +\mu _i\mu _j=E\left( \varepsilon _i\varepsilon _j \right) -\mu _j\mu _i-\mu _i\mu _j+\mu _i\mu _j
\\
&=E\left( \varepsilon _i\varepsilon _j \right) -\mu _i\mu _j.
\end{align*}
进而
\begin{align*}
E\left( \varepsilon _i\varepsilon _j \right) =\mathrm{Cov}\left( \varepsilon _i,\varepsilon _j \right) +\mu _i\mu _j.
\end{align*}
将上式代入\eqref{eq::HJIOJW392234--32}式得
\begin{align*}
E(\varepsilon^T A\varepsilon )&=\sum_{i,j=1}^n{a_{ij}\Sigma _{ij}}+\sum_{i,j=1}^n{a_{ij}\mu _i\mu _j}=\sum_{i,j=1}^n{a_{ij}\Sigma _{ji}}+\mu ' A\mu 
\\
&=\sum_{i=1}^n{(A\Sigma )_{ii}}+\mu' A\mu =\mathrm{tr(}A\Sigma )+\mu' A\mu .
\end{align*}
特别地,若$\mu=0$，$\Sigma=\sigma^2 I_n$，代入上式得$E(\varepsilon^TA\varepsilon)=\sigma^2\operatorname{tr}(A)$。
\end{proof}

\begin{theorem}[OLS误差方差的无偏估计]
考虑线性回归模型
\[
Y = X\beta + \varepsilon,
\]
其中：
\begin{itemize}
\item $Y$是$n$维观测向量，
\item $X$ 是 $n \times (K+1)$列满秩设计矩阵$(K+1<n)$，
\item $\beta$ 是 $(K+1)$维未知参数向量，
\item $\varepsilon$ 是 $n$维随机误差向量。
\end{itemize}

假设：
\begin{enumerate}[(1)]
\item $E(\varepsilon)=0$，
\item $\operatorname{Cov}(\varepsilon)=E(\varepsilon\varepsilon^T)=\sigma^2 I_n$，其中 $\sigma^2=\text{Var}(\varepsilon_i)>0(i=1,2,\cdots,n)$是未知常数.
\end{enumerate}

定义：
\begin{itemize}
\item OLS 估计量：$\hat{\beta}=(X^TX)^{-1}X^TY$，
\item 残差向量：$e=Y-X\hat{\beta}$，
\item 残差平方和：$e^Te$，
\item 误差方差估计量：$s_e^2=\dfrac{e^Te}{n-K-1}$。
\end{itemize}
则 $s_e^2$ 是 $\sigma^2$ 的无偏估计，即
\[
E(s_e^2)=\sigma^2.
\]
\end{theorem}
\begin{proof}

由题设可知
\[
\begin{aligned}
e &= Y - X\hat{\beta} = (X\beta+\varepsilon) - X[(X^TX)^{-1}X^T(X\beta+\varepsilon)] \\
&= X\beta+\varepsilon - X\beta - X(X^TX)^{-1}X^T\varepsilon \\
&= \bigl[I_n - X(X^TX)^{-1}X^T\bigr]\varepsilon.
\end{aligned}
\]
记 $M = I_n - X(X^TX)^{-1}X^T$，则 $e=M\varepsilon$。显然$M$是对称幂等阵,从而
\[
e^Te = (M\varepsilon)^T(M\varepsilon)=\varepsilon^T M'M \varepsilon = \varepsilon^T M \varepsilon,
\]
由假设知 $\varepsilon$ 满足 $E(\varepsilon)=0$，$\operatorname{Cov}(\varepsilon)=\sigma^2 I_n$，则由\reflem{lemma:二次型期望公式}得
\begin{align*}
E\left( e^Te \right) &=\sigma ^2\mathrm{tr}\left( M \right) =\sigma ^2\mathrm{tr}\bigl( I_n-X\left( X^TX \right) ^{-1}X^T \bigr) 
\\
&=\sigma ^2\mathrm{tr}\left( I_n \right) -\sigma ^2\mathrm{tr}\bigl( X\left( X^TX \right) ^{-1}X^T \bigr) =\sigma ^2n-\sigma ^2\mathrm{tr}\bigl( \left( X^TX \right) ^{-1}X^TX \bigr) 
\\
&=\sigma ^2n-\sigma ^2\mathrm{tr}\left( I_{K+1} \right) =\sigma ^2\left( n-K-1 \right) .
\end{align*}
因此
\[
E(e^Te)=\sigma^2\bigl(n-(K+1)\bigr)=\sigma^2(n-K-1).
\]
故
\[
E(s_e^2)=E\left(\frac{e^Te}{n-K-1}\right)=\frac{1}{n-K-1}E(e^Te)=\frac{1}{n-K-1}\sigma^2(n-K-1)=\sigma^2.
\]
\end{proof}




















\end{document}